\chapter{Saisie des transactions et enregistrement}\label{ch:pl:trade-capture-and-booking}

Un swap de 100~millions de dollars a été modifié pour atteindre 150~millions de dollars le lundi, a fait l'objet d'une novation au profit d'un autre dealer le mercredi et a été partiellement résilié le vendredi
matin. Le vendredi après-midi, le système de risque détenait la version~3 (150~millions de dollars, le nouveau dealer), le système de confirmation la version~2 (150~millions de dollars,
l'ancien dealer) et le système de règlement-livraison la version~4 (90~millions de dollars, le nouveau dealer). Chacun avait copié la transaction lorsqu'un événement l'intéressait,
mais aucun n'avait suivi ce qui lui était arrivé. Ce chapitre traite une transaction comme le résultat de ses événements~: saisie une seule fois
dans un modèle canonique, enregistrée selon des règles, allouée selon une règle énoncée, et modifiée uniquement par des événements que lisent tous les systèmes en aval.

\section{Ce qu'est une transaction, comme donnée}

\begin{definition}[Modèle canonique de transaction]\label{def:pl:trade-capture-and-booking:canonical}
Un \emph{modèle canonique de transaction}\index{modèle canonique de transaction} est la représentation unique qu'a la société d'une transaction, partagée par tous les systèmes qui
la traitent~: un en-tête (identifiants, parties, book, dates, statut, version) et l'économie (l'instrument, la quantité et le prix),
assortis de règles que toute transaction valide satisfait.
\end{definition}

\begin{definition}[Identifiant unique de transaction (UTI)]\label{def:pl:trade-capture-and-booking:uti}
Un \emph{identifiant unique de transaction}\index{identifiant unique de transaction} (UTI) identifie une transaction pour l'ensemble des parties et
des référentiels qui la déclarent~: selon les recommandations techniques du CPMI--IOSCO, l'identifiant d'entité juridique (LEI) (livre~2, chapitre~28) de l'entité qui
la génère, suivi d'une valeur unique pour cette entité, de 52 caractères au plus, composée uniquement de lettres majuscules et de chiffres.
\end{definition}

Le modèle du chapitre (\texttt{firm.tradecap}) conserve l'instrument sous la forme dans laquelle la bibliothèque de valorisation du livre~5 le sérialise, de sorte qu'une transaction saisie
soit valorisée, gérée au regard du risque et déclarée à partir du même \omterm{def:pl:trade-capture-and-booking:capture}{enregistrement}~; le validateur contrôle le format des identifiants, que l'acheteur et le vendeur diffèrent, que
le règlement-livraison ne précède pas la transaction, et que la quantité est positive. Les modèles propres au secteur --- FpML pour les messages, le Common Domain
Model de l'ISDA pour les événements --- reprennent la même idée, à une échelle bien plus grande.

\section{Saisie~: de l'exécution à la transaction enregistrée}

\begin{definition}[Saisie des transactions, enregistrement]\label{def:pl:trade-capture-and-booking:capture}
La \emph{saisie des transactions}\index{saisie des transactions} (trade capture) transforme une exécution --- un rapport d'exécution de bourse, un ticket OTC, une transaction conclue par téléphone --- en transaction
dans le modèle canonique, une seule fois et le plus tôt possible. L'\emph{enregistrement}\index{enregistrement} (booking) affecte la transaction saisie à un book, l'unité qui porte
son risque et son P\&L.
\end{definition}

\begin{definition}[Modèle d'enregistrement]\label{def:pl:trade-capture-and-booking:model}
Un \emph{modèle d'enregistrement}\index{modèle d'enregistrement} (booking model) est l'ensemble des règles qui décide où une transaction est enregistrée --- selon le type d'instrument, le desk, l'entité
juridique et la stratégie --- de sorte que le book découle des attributs de la transaction, et non d'un champ saisi à la main.
\end{definition}

La partie boursière du chapitre est un bloc~: 9\,900 actions achetées sur le simulateur du livre~10 par un algorithme d'exécution envoyant un ordre qui prend le
cours vendeur toutes les 20~secondes, exécuté en 57 exécutions sur trois minutes, à un prix moyen de 100.0783~dollars. Chaque rapport d'exécution devient une transaction dotée de sa
propre UTI, achetée par la société auprès de la chambre de compensation. La partie OTC porte sur 200 swaps venant de six dealers, saisis à partir de tickets le lundi matin.

\begin{omfigure}
\begin{tikzpicture}
\begin{axis}[width=11.5cm,height=5cm,xlabel={secondes après l'ouverture},ylabel={prix (\$)},grid=major,grid style={gray!20},
  tick label style={font=\small},label style={font=\small},table/col sep=comma,
  scatter,only marks,point meta={\thisrow{qty}},scatter/use mapped color={draw=omDef,fill=omDef!40},
  visualization depends on={sqrt(\thisrow{qty})/6 \as \msize},scatter/@pre marker code/.append style={/tikz/mark size=\msize}]
\addplot table[x=seconds,y=price]{figdata/platforms/21-trade-capture-and-booking/block.csv};
\draw[omThm,dashed] (axis cs:0,100.0783) -- (axis cs:220,100.0783) node[pos=0.02,above right,font=\scriptsize]{moyenne 100.0783~dollars};
\end{axis}
\end{tikzpicture}

\omcaption{Les 57 exécutions du bloc de 9\,900 actions sur le simulateur, selon le temps et le prix (aire des marqueurs proportionnelle à la taille)~: dix ordres enfants prenant le
cours vendeur toutes les 20~secondes, chacun étant exécuté en plusieurs exécutions au fur et à mesure qu'il parcourt le book. Chaque exécution est saisie comme une transaction~; chaque fonds est alloué
au prix moyen. Données~: \texttt{fig\_tradecap.py}.}\label{fig:pl:trade-capture-and-booking:block}
\end{omfigure}

\section{Le cycle de vie en événements}

\begin{definition}[Événement de cycle de vie de la transaction, modification de transaction]\label{def:pl:trade-capture-and-booking:event}
Un \emph{événement de cycle de vie de la transaction}\index{evenement@événement de cycle de vie de la transaction} est tout ce qui arrive à une transaction après son exécution~: une modification, une
novation (livre~1, chapitre~5), une résiliation partielle, un exercice, une allocation, une annulation. Une \emph{modification de transaction}\index{modification de transaction}
modifie une condition de la transaction (notionnel, taux, dates) d'un commun accord~; elle est enregistrée comme une nouvelle version, jamais en écrasant l'ancienne.
\end{definition}

Le Common Domain Model de l'ISDA décrit les événements de cycle de vie dans les mêmes termes~: «~un événement de cycle de vie décrit une transition d'état,~» avec un
état antérieur et un état postérieur, composée à partir de quelques instructions élémentaires --- exécution, changement de quantité, changement de conditions, changement de partie, exercice,
transfert et une poignée d'autres. \texttt{firm.tradecap} conserve les événements et calcule la transaction~: un magasin applique les événements dans l'ordre chronologique et
reconstitue n'importe quelle version à n'importe quel instant (\cref{lst:pl:trade-capture-and-booking:store}). L'event sourcing du chapitre~24 est le
schéma général~; ici, c'est la seule façon de savoir, le vendredi, ce qu'était la transaction le mercredi.

\omcode{code/firm/tradecap/firm_tradecap.py}{70}{98}{Le magasin de transactions~: événements ajoutés, jamais modifiés~; une transaction à un instant quelconque est le rejeu
de ses événements jusqu'à cet instant.}\label{lst:pl:trade-capture-and-booking:store}

\omcode{code/firm/tradecap/firm_tradecap.py}{101}{124}{Chaque type d'événement comme une transition d'une version de la transaction à la
version suivante.}\label{lst:pl:trade-capture-and-booking:step}

\begin{omfigure}
\begin{tikzpicture}[font=\scriptsize,>=Stealth,ver/.style={draw,rounded corners=2pt,align=center,minimum height=1.0cm,minimum width=2.2cm}]
\node[ver,fill=gray!10] (v1) at (0,0) {v1~: 100~M\$\\ dealer 1};
\node[ver,fill=gray!10] (v2) at (3.4,0) {v2~: 150~M\$\\ dealer 1};
\node[ver,fill=gray!10] (v3) at (6.8,0) {v3~: 150~M\$\\ dealer 4};
\node[ver,fill=omDef!15] (v4) at (10.2,0) {v4~: 90~M\$\\ dealer 4};
\draw[->] (v1) -- node[above=4pt,align=center]{modif.\\ lun.} (v2);
\draw[->] (v2) -- node[above=4pt,align=center]{novation\\ mer.} (v3);
\draw[->] (v3) -- node[above=4pt,align=center]{réduction de 40~\%\\ ven.} (v4);
\node[omThm,align=center] at (3.4,-1.6) {confirmation\\ (copies à la création,\\ à la modification)};
\node[omThm,align=center] at (6.8,-1.6) {risque\\ (copies chaque soir)};
\node[omThm,align=center] at (10.2,-1.6) {règlement-livraison\\ (copies à la création,\\ à la résiliation)};
\draw[->,omThm] (3.4,-0.95) -- (v2.south); \draw[->,omThm] (6.8,-0.95) -- (v3.south); \draw[->,omThm] (10.2,-0.95) -- (v4.south);
\end{tikzpicture}

\omcaption{Le swap de l'accroche au fil de la semaine et ce que trois systèmes fondés sur la copie détenaient le vendredi à 15~h~00. Chacun a copié la transaction sur les événements
qui le concernent, ou selon son propre calendrier. Le règlement-livraison ne détient la version~4 que parce que le dernier événement, une résiliation, fait partie de ceux sur lesquels il copie~; seul un
système qui suit chaque événement détient la version courante, quel que soit le dernier événement.}\label{fig:pl:trade-capture-and-booking:versions}
\end{omfigure}

\section{Trois versions le vendredi}

Le cycle de vie de la semaine est simulé~: 69 modifications, 26 novations et 41 résiliations partielles parmi les 200 swaps entre le lundi 10~h~00 et le vendredi
à 14~h~00, dont le swap de l'accroche. Trois systèmes en aval conservent des copies~: le système de risque copie chaque transaction à 18~h~00, chaque soir~; la confirmation reçoit une copie sur
les événements qu'elle confirme (nouvelles transactions et modifications)~; le règlement-livraison sur les événements qui mettent en jeu des flux de trésorerie (nouvelles transactions et résiliations). Le vendredi à 15~h~00
leurs copies sont comparées à la transaction issue de l'event sourcing.

\omcode{code/platforms/21-trade-capture-and-booking/python/pl_tradecap.py}{118}{144}{Les trois politiques de copie et la comparaison avec la
transaction issue de l'event sourcing, le vendredi après-midi.}\label{lst:pl:trade-capture-and-booking:downstream}

\begin{table}[ht]
\centering\small
\begin{tabular}{lrr}
\toprule
système & copie fautive & événements fautifs \\
\midrule
risque (copie nocturne) & 25 & 0 \\
confirmation (copie à la création et à la modification) & 60 & 0 \\
règlement-livraison (copie à la création et à la résiliation) & 71 & 0 \\
\midrule
transactions pour lesquelles les trois divergent & 109 sur 200 & 0 \\
\bottomrule
\end{tabular}
\caption{Swaps dont la copie en aval diverge de la transaction issue de l'event sourcing le vendredi à 15~h~00, après une semaine
d'événements de cycle de vie.}\label{tab:pl:trade-capture-and-booking:downstream}
\end{table}

\begin{omfigure}
\begin{tikzpicture}
\begin{axis}[width=11.5cm,height=5cm,xbar,bar width=8pt,xmin=0,xmax=130,xlabel={swaps fautifs le vendredi à 15~h~00 (sur 200)},
  symbolic y coords={any disagreement,settlement,confirmation,risk},ytick=data,yticklabels={risque,confirmation,règlement,tout désaccord},enlarge y limits=0.2,
  tick label style={font=\small},label style={font=\small},grid=major,grid style={gray!20},table/col sep=comma,
  nodes near coords,nodes near coords style={font=\scriptsize,anchor=west},legend style={font=\footnotesize,at={(0.5,1.03)},
  anchor=south,legend columns=2}]
\addplot[fill=omThm!45,draw=omThm] table[y=system,x=copy_wrong]{figdata/platforms/21-trade-capture-and-booking/downstream.csv};
\addlegendentry{copie}
\addplot[fill=omDef!45,draw=omDef] table[y=system,x=events_wrong]{figdata/platforms/21-trade-capture-and-booking/downstream.csv};
\addlegendentry{suivi des événements}
\end{axis}
\end{tikzpicture}

\omcaption{Systèmes en aval comparés aux swaps issus de l'event sourcing le vendredi après-midi, après une semaine de modifications, de novations et de résiliations
partielles~: les copies sont fausses dans 25 à 71 transactions et divergent sur 109~; les systèmes qui suivent les événements ne se trompent jamais. Données~:
\texttt{fig\_tradecap.py}.}\label{fig:pl:trade-capture-and-booking:downstream}
\end{omfigure}

La copie échoue de façon différente dans chaque système (\cref{tab:pl:trade-capture-and-booking:downstream})~: le risque n'est fautif que pour les événements du jour
(25 transactions), la confirmation pour chaque novation et chaque résiliation qu'elle n'a jamais reçue (60), le règlement-livraison pour chaque modification et chaque novation (71), et
sur 109 des 200 swaps, au moins deux systèmes divergent. Suivant les événements, les trois détiennent la même transaction, celle que reconstruit le magasin.

\section{Allocations et traitement de bout en bout}

\begin{definition}[Allocation de bloc]\label{def:pl:trade-capture-and-booking:allocation}
Une \emph{allocation de bloc}\index{allocation de bloc} répartit une transaction exécutée pour plusieurs comptes à la fois --- un bloc pour plusieurs fonds --- entre
eux après l'exécution, au prix moyen du bloc, selon une règle fixée avant la transaction, qui décide aussi du sort des unités qui ne se divisent
pas exactement.
\end{definition}

\omcode{code/firm/tradecap/firm_tradecap.py}{144}{166}{Allocation par lots entiers au prix moyen, avec trois règles pour les lots qui ne
se divisent pas.}\label{lst:pl:trade-capture-and-booking:allocate}

Le bloc de 9\,900 actions (99 lots de 100) est alloué à trois fonds dont les cibles sont de 50~\%, 30~\% et 20~\%, et dont les parts exactes sont 4\,950, 2\,970
et 1\,980. La règle du plus fort reste donne 4\,900, 3\,000 et 2\,000~: les 99 lots sont alloués, et aucun fonds ne s'écarte de plus de 50 actions de sa cible. Attribuer
les lots irréguliers aux plus gros fonds donne 5\,000, 3\,000 et 1\,900, soit 80 actions d'écart par rapport à la cible pour le plus petit. Arrondir chaque fonds séparément donne
5\,000, 3\,000 et 2\,000 --- 10\,000 actions pour un bloc de 9\,900~: une sur-allocation d'un lot que quelqu'un doit acheter ou restituer.
(\cref{tab:pl:trade-capture-and-booking:allocation}). Chaque fonds paie le prix moyen du bloc, 100.0783~\$, quelles que soient les exécutions qui l'ont produit.

\begin{table}[ht]
\centering\small
\begin{tabular}{lrrrrr}
\toprule
règle & fonds A & fonds B & fonds C & reliquat & écart maximal \\
\midrule
plus fort reste & 4\,900 & 3\,000 & 2\,000 & 0 & 50 \\
lots irréguliers aux plus gros fonds & 5\,000 & 3\,000 & 1\,900 & 0 & 80 \\
chaque fonds arrondi & 5\,000 & 3\,000 & 2\,000 & $-100$ & 50 \\
\bottomrule
\end{tabular}
\caption{Un bloc de 9\,900 actions alloué à trois fonds (cibles 50~\%, 30~\%, 20~\%) par lots de 100 selon trois règles~: actions par fonds, le
bloc moins l'allocation, et l'écart maximal à la part exacte d'un fonds.}\label{tab:pl:trade-capture-and-booking:allocation}
\end{table}

\begin{definition}[Traitement de bout en bout]\label{def:pl:trade-capture-and-booking:stp}
Le \emph{traitement de bout en bout}\index{traitement de bout en bout} (straight-through processing) mène une transaction de l'exécution au règlement --- saisie, \omterm{def:pl:trade-capture-and-booking:capture}{enregistrement},
allocation, confirmation, compensation, règlement --- sans ressaisie ni intervention manuelle, chaque étape lisant la sortie de la précédente dans le
modèle canonique.
\end{definition}

Le \omterm{def:pl:trade-capture-and-booking:stp}{traitement de bout en bout} est ce que permettent le modèle canonique et le flux d'événements~: chaque étape manuelle est une copie, et chaque copie est une version qui peut
diverger. Le taux de transactions qui passent sans intervention humaine est la mesure opérationnelle de la \omterm{def:pl:trade-capture-and-booking:capture}{saisie des transactions} d'une plateforme (le chapitre~22 mesure
les ruptures lorsque ce n'est pas le cas).

\begin{dated}{2026-09}{Identifiants et modèles d'événements}\label{dat:pl:trade-capture-and-booking:standards}
Le CPMI et l'IOSCO ont publié leurs recommandations techniques sur l'harmonisation de l'UTI le 28 février 2017~: le LEI de l'entité qui génère l'identifiant, concaténé
à une valeur unique pour cette entité, sans chiffre de contrôle, au plus 52 caractères parmi A--Z et 0--9. Le Common Domain Model de l'ISDA, maintenu en source ouverte,
modélise les événements de cycle de vie comme des transitions d'état composées d'instructions élémentaires (exécution, changement de quantité, changement de conditions, changement de contrepartie,
exercice et autres).
\end{dated}

\section{Tutoriel~: un bloc, deux cents swaps, une semaine}\label{tut:pl:trade-capture-and-booking:tut}

\begin{tutorial}
\textbf{Objectif.} Saisir les exécutions et les tickets dans un seul modèle, allouer un bloc, faire tourner une semaine d'événements de cycle de vie, et comparer la copie avec
le suivi des événements. \textbf{État final~:} \cref{tab:pl:trade-capture-and-booking:downstream,tab:pl:trade-capture-and-booking:allocation}.
\begin{tutsteps}
\item \textbf{Saisie}~: \texttt{pl\_tradecap.block\_fills()} et \texttt{capture\_block}~; valider chaque transaction.
\item \textbf{Allocation}~: \texttt{allocations(fills)} selon les trois règles.
\item \textbf{Cycle de vie}~: \texttt{week()}~; \texttt{store.versions} et \texttt{store.as\_of}.
\item \textbf{Aval}~: \texttt{downstream(store, events)}.
\end{tutsteps}
\textbf{Pour aller plus loin.} Ajouter un événement d'exercice qui transforme une swaption en swap et le suivre en aval~; allouer selon les liquidités des
fonds au lieu de cibles fixes.
\end{tutorial}

\section{Construction~: saisie des transactions}\label{bld:pl:trade-capture-and-booking:tradecap}

\begin{build}
\bfield{Rôle} Un seul \omterm{def:pl:trade-capture-and-booking:capture}{enregistrement} de transaction pour tous les systèmes, modifié uniquement par des événements, de la saisie au règlement.
\bfield{Interface} \texttt{uti}, \texttt{validate}, \texttt{Event}, \texttt{TradeStore} (\texttt{apply}, \texttt{as\_of},
\texttt{versions}), \texttt{from\_execution}, \texttt{from\_ticket}, \texttt{allocate}, \texttt{BookingModel}.
\bfield{Règles} Les transactions ne sont jamais modifiées directement, seulement par des événements~; chaque transaction est validée~; les identifiants suivent les formats de l'UTI et du LEI~; les règles d'allocation sont
fixées avant la transaction et allouent le bloc entier~; les livres proviennent des règles.
\bfield{Tests d'acceptation} \texttt{code/firm/tradecap/tests/}~: format de l'UTI et erreurs de validation~; les versions et les états d'un cycle de vie à tout instant
, y compris la résiliation à zéro et un événement sur une transaction inconnue~; l'exercice et la saisie à partir d'un rapport d'exécution~; les trois règles
d'allocation~; le \omterm{def:pl:trade-capture-and-booking:model}{modèle d'enregistrement}.
\bfield{Pour aller plus loin} Messages en FpML ou événements au format JSON du CDM~; allocations par liquidités avec parts fractionnaires lorsque cela est permis~; un rapport sur le taux de
\omterm{def:pl:trade-capture-and-booking:stp}{traitement de bout en bout}.
\end{build}

\begin{omsources}
\item CPMI et IOSCO, \emph{Harmonisation of the Unique Transaction Identifier: Technical Guidance}, février 2017.
\item Documentation du Common Domain Model de l'ISDA, «~Event Model~».
\item Livre~1, chapitre~5, et livre~2, chapitre~28 (novation, LEI)~; livre~10, chapitre~26 (rapports d'exécution).
\end{omsources}

\section{Exercices}

\begin{exercise}[$\star$]\label{exo:pl:trade-capture-and-booking:1}
Un bloc de 1\,000 actions est réparti entre trois fonds à parts égales, par lots de 100. Que donne chaque règle~?
\end{exercise}

\begin{exercise}[$\star$]\label{exo:pl:trade-capture-and-booking:2}
Pourquoi le système de risque se trompe-t-il sur moins de transactions que les deux autres~?
\end{exercise}

\begin{exercise}[$\star$]\label{exo:pl:trade-capture-and-booking:3}
À quoi ressemble un UTI généré par la société pour sa 7\up{e} transaction~?
\end{exercise}

\begin{exercise}[$\star\star$]\label{exo:pl:trade-capture-and-booking:4}
Pourquoi la règle d'allocation doit-elle être fixée avant la transaction, et que se passe-t-il si un trader la choisit après avoir vu les exécutions~?
\end{exercise}

\begin{exercise}[$\star\star$]\label{exo:pl:trade-capture-and-booking:5}
Le système de confirmation se trompe sur 60 transactions. Quels événements a-t-il manqués, et à quoi aurait-il dû s'abonner~?
\end{exercise}

\begin{exercise}[$\star\star$]\label{exo:pl:trade-capture-and-booking:6}
Une transaction a été enregistrée dans le mauvais livre, puis déplacée le lendemain. Comment le déplacement doit-il être consigné, et pourquoi~?
\end{exercise}

\begin{exercise}[$\star\star\star$]\label{exo:pl:trade-capture-and-booking:7}
\emph{Programmation.} Relancer \texttt{downstream} avec le système de risque qui copie aussi à 12~h~00, en plus de 18~h~00. Sur combien de transactions se trompe-t-il~?
\end{exercise}

\begin{exercise}[$\star\star\star$]\label{exo:pl:trade-capture-and-booking:8}
\emph{Trouver la faille.} «~Chaque système possède la transaction, donc ils doivent concorder~: chacun la lit dans le système de trading lorsqu'il en a besoin.~»
\end{exercise}

\section{Problème~: Trois versions le vendredi}

\begin{problem}[{Problème du week-end --- copies contre événements}]\label{pb:pl:trade-capture-and-booking:1}
Le bloc, les 200 swaps et la semaine du chapitre.

\medskip\noindent\textbf{Partie I --- Le modèle.}
\begin{enumerate}
\item Que contient le modèle canonique, et pourquoi l'instrument tel que le sérialise la bibliothèque de pricing~?
\item Comment construit-on un UTI~?
\item Que vérifie le validateur~?
\item Que décide un \omterm{def:pl:trade-capture-and-booking:model}{modèle d'enregistrement}~?
\item Pourquoi une transaction n'est-elle jamais modifiée~?
\end{enumerate}

\medskip\noindent\textbf{Partie II --- La semaine.}
\begin{enumerate}[resume]
\item Quels événements ont touché les swaps~?
\item Par quoi est passé le swap du hook, et que détenait chaque système le vendredi~?
\item Que manque-t-il à chaque politique de copie~?
\item Sur combien de transactions chaque système s'est-il trompé~?
\item Sur combien de transactions les systèmes ont-ils divergé, et combien sous le régime des événements~?
\end{enumerate}

\medskip\noindent\textbf{Partie III --- Le bloc.}
\begin{enumerate}[resume]
\item Comment le bloc a-t-il été exécuté, et à quel prix moyen~?
\item Quelles sont les parts exactes des trois fonds~?
\item Que répartit chaque règle, et avec quel reliquat~?
\item Quel est l'écart maximal par rapport à la cible pour chaque règle~?
\item Quel prix chaque fonds paie-t-il~?
\end{enumerate}

\medskip\noindent\textbf{Partie IV --- Le verdict.}
\begin{enumerate}[resume]
\item Énoncer le \emph{résultat nommé}~: le nombre de divergences en aval après une semaine d'événements de cycle de vie, selon que la diffusion se fait par copies
ou par événements, et les reliquats d'allocation d'un bloc réparti entre trois fonds selon trois règles d'arrondi.
\item Quelle règle d'allocation faudrait-il adopter, et pourquoi~?
\item Que faudrait-il exiger d'un nouveau système en aval~?
\item Que mesure le \omterm{def:pl:trade-capture-and-booking:stp}{traitement de bout en bout}~?
\item En une phrase~: qu'est-ce qu'une transaction~?
\end{enumerate}
\end{problem}

\section{Questions d'entretien}

\begin{interviewq}[$\star$ \iqroles{développeur}]\label{iq:pl:trade-capture-and-booking:1}
Trois systèmes affichent trois versions de la même transaction. Comment découvrez-vous laquelle est juste, et comment empêchez-vous que cela se reproduise~?
\end{interviewq}

\begin{interviewq}[$\star\star$ \iqroles{développeur}]\label{iq:pl:trade-capture-and-booking:2}
Comment représenteriez-vous une transaction pour que le pricing, le risque, la confirmation et le règlement puissent tous l'utiliser~?
\end{interviewq}

\begin{interviewq}[$\star\star$ \iqroles{développeur}]\label{iq:pl:trade-capture-and-booking:3}
Un bloc est exécuté à vingt prix différents pour cinq fonds. Comment l'allouez-vous~?
\end{interviewq}

\begin{interviewq}[$\star\star$ \iqroles{développeur}]\label{iq:pl:trade-capture-and-booking:4}
Qu'est-ce qu'une novation, et que change-t-elle dans les systèmes qui détiennent la transaction~?
\end{interviewq}

\begin{interviewq}[$\star\star\star$ \iqroles{développeur}]\label{iq:pl:trade-capture-and-booking:5}
Concevez la couche de saisie et d'\omterm{def:pl:trade-capture-and-booking:capture}{enregistrement} des transactions d'une société de trading multi-actifs.
\end{interviewq}
